% Kartikey Pathak --- one-page generalist resume (the default one to send)
% build:  pdflatex -interaction=nonstopmode resume.tex
\documentclass[a4paper,9pt]{extarticle}
\usepackage[left=1.2cm,right=1.2cm,top=0.8cm,bottom=0.65cm]{geometry}

\newcommand{\secabove}{0.42em}
\newcommand{\secbelow}{0.26em}
\newcommand{\pointsep}{0.03em}
\newcommand{\pointafter}{0.12em}

% Shared preamble for both resume variants.
% Set \DENSE to 1 before % Shared preamble for both resume variants.
% Set \DENSE to 1 before % Shared preamble for both resume variants.
% Set \DENSE to 1 before \input{preamble} for the tighter one-page build.

\usepackage[T1]{fontenc}
\usepackage[utf8]{inputenc}
\usepackage{charter}
\usepackage{titlesec}
\usepackage{enumitem}
\usepackage{xcolor}
\usepackage[hidelinks]{hyperref}
\usepackage{tabularx}

\definecolor{ink}{HTML}{111111}
\definecolor{link}{HTML}{1A4B8C}
\color{ink}
\hypersetup{colorlinks=true,urlcolor=link,linkcolor=link}

\pagestyle{empty}
\setlength{\parindent}{0pt}
\raggedright

% ---- section headings -------------------------------------------------
\titleformat{\section}
  {\normalfont\large\bfseries\scshape}{}{0em}{}[\vspace{-0.72em}\rule{\linewidth}{0.6pt}]
\titlespacing{\section}{0pt}{\secabove}{\secbelow}

% ---- helpers ----------------------------------------------------------
% \entry{organisation}{dates}{role}{location}
\newcommand{\entry}[4]{%
  \begin{tabularx}{\linewidth}{@{}Xr@{}}
    \textbf{#1} & #2 \\
    \textit{#3} & \textit{#4} \\
  \end{tabularx}%
  \vspace{-0.15em}%
}
% \pentry{project}{stack}{dates}
\newcommand{\pentry}[3]{%
  \begin{tabularx}{\linewidth}{@{}Xr@{}}
    \textbf{#1} \,\textbar\, \textit{#2} & #3 \\
  \end{tabularx}%
  \vspace{-0.15em}%
}
\newenvironment{points}
  {\begin{itemize}[leftmargin=1.1em,itemsep=\pointsep,topsep=0.18em,parsep=0pt,label=--]}
  {\end{itemize}\vspace{\pointafter}}
\newcommand{\skillrow}[2]{\textbf{#1:} #2\\[0.12em]}

% ---- the header, identical in both variants ---------------------------
\newcommand{\resumeheader}{%
\begin{center}
  {\LARGE\bfseries Kartikey Pathak}\\[0.35em]
  {\small I build machines end to end --- mechanical, PCB, firmware, control, and the app on top.}\\[0.4em]
  {\small
    +91 7506048765 \,\textbar\,
    \href{mailto:kartikeypathak.imp@gmail.com}{kartikeypathak.imp@gmail.com} \,\textbar\,
    \href{https://spongefal.vercel.app}{spongefal.vercel.app} \,\textbar\,
    \href{https://github.com/NoobMaster-version}{github.com/NoobMaster-version} \,\textbar\,
    \href{https://www.linkedin.com/}{LinkedIn}
  }
\end{center}
}
 for the tighter one-page build.

\usepackage[T1]{fontenc}
\usepackage[utf8]{inputenc}
\usepackage{charter}
\usepackage{titlesec}
\usepackage{enumitem}
\usepackage{xcolor}
\usepackage[hidelinks]{hyperref}
\usepackage{tabularx}

\definecolor{ink}{HTML}{111111}
\definecolor{link}{HTML}{1A4B8C}
\color{ink}
\hypersetup{colorlinks=true,urlcolor=link,linkcolor=link}

\pagestyle{empty}
\setlength{\parindent}{0pt}
\raggedright

% ---- section headings -------------------------------------------------
\titleformat{\section}
  {\normalfont\large\bfseries\scshape}{}{0em}{}[\vspace{-0.72em}\rule{\linewidth}{0.6pt}]
\titlespacing{\section}{0pt}{\secabove}{\secbelow}

% ---- helpers ----------------------------------------------------------
% \entry{organisation}{dates}{role}{location}
\newcommand{\entry}[4]{%
  \begin{tabularx}{\linewidth}{@{}Xr@{}}
    \textbf{#1} & #2 \\
    \textit{#3} & \textit{#4} \\
  \end{tabularx}%
  \vspace{-0.15em}%
}
% \pentry{project}{stack}{dates}
\newcommand{\pentry}[3]{%
  \begin{tabularx}{\linewidth}{@{}Xr@{}}
    \textbf{#1} \,\textbar\, \textit{#2} & #3 \\
  \end{tabularx}%
  \vspace{-0.15em}%
}
\newenvironment{points}
  {\begin{itemize}[leftmargin=1.1em,itemsep=\pointsep,topsep=0.18em,parsep=0pt,label=--]}
  {\end{itemize}\vspace{\pointafter}}
\newcommand{\skillrow}[2]{\textbf{#1:} #2\\[0.12em]}

% ---- the header, identical in both variants ---------------------------
\newcommand{\resumeheader}{%
\begin{center}
  {\LARGE\bfseries Kartikey Pathak}\\[0.35em]
  {\small I build machines end to end --- mechanical, PCB, firmware, control, and the app on top.}\\[0.4em]
  {\small
    +91 7506048765 \,\textbar\,
    \href{mailto:kartikeypathak.imp@gmail.com}{kartikeypathak.imp@gmail.com} \,\textbar\,
    \href{https://spongefal.vercel.app}{spongefal.vercel.app} \,\textbar\,
    \href{https://github.com/NoobMaster-version}{github.com/NoobMaster-version} \,\textbar\,
    \href{https://www.linkedin.com/}{LinkedIn}
  }
\end{center}
}
 for the tighter one-page build.

\usepackage[T1]{fontenc}
\usepackage[utf8]{inputenc}
\usepackage{charter}
\usepackage{titlesec}
\usepackage{enumitem}
\usepackage{xcolor}
\usepackage[hidelinks]{hyperref}
\usepackage{tabularx}

\definecolor{ink}{HTML}{111111}
\definecolor{link}{HTML}{1A4B8C}
\color{ink}
\hypersetup{colorlinks=true,urlcolor=link,linkcolor=link}

\pagestyle{empty}
\setlength{\parindent}{0pt}
\raggedright

% ---- section headings -------------------------------------------------
\titleformat{\section}
  {\normalfont\large\bfseries\scshape}{}{0em}{}[\vspace{-0.72em}\rule{\linewidth}{0.6pt}]
\titlespacing{\section}{0pt}{\secabove}{\secbelow}

% ---- helpers ----------------------------------------------------------
% \entry{organisation}{dates}{role}{location}
\newcommand{\entry}[4]{%
  \begin{tabularx}{\linewidth}{@{}Xr@{}}
    \textbf{#1} & #2 \\
    \textit{#3} & \textit{#4} \\
  \end{tabularx}%
  \vspace{-0.15em}%
}
% \pentry{project}{stack}{dates}
\newcommand{\pentry}[3]{%
  \begin{tabularx}{\linewidth}{@{}Xr@{}}
    \textbf{#1} \,\textbar\, \textit{#2} & #3 \\
  \end{tabularx}%
  \vspace{-0.15em}%
}
\newenvironment{points}
  {\begin{itemize}[leftmargin=1.1em,itemsep=\pointsep,topsep=0.18em,parsep=0pt,label=--]}
  {\end{itemize}\vspace{\pointafter}}
\newcommand{\skillrow}[2]{\textbf{#1:} #2\\[0.12em]}

% ---- the header, identical in both variants ---------------------------
\newcommand{\resumeheader}{%
\begin{center}
  {\LARGE\bfseries Kartikey Pathak}\\[0.35em]
  {\small I build machines end to end --- mechanical, PCB, firmware, control, and the app on top.}\\[0.4em]
  {\small
    +91 7506048765 \,\textbar\,
    \href{mailto:kartikeypathak.imp@gmail.com}{kartikeypathak.imp@gmail.com} \,\textbar\,
    \href{https://spongefal.vercel.app}{spongefal.vercel.app} \,\textbar\,
    \href{https://github.com/NoobMaster-version}{github.com/NoobMaster-version} \,\textbar\,
    \href{https://www.linkedin.com/}{LinkedIn}
  }
\end{center}
}


\begin{document}
\resumeheader

% =======================================================================
\section{Education}
\entry{Veermata Jijabai Technological Institute (VJTI)}{2023 -- 2027}
      {B.Tech, Electrical Engineering --- Minor in Robotics (CGPA 8.04)}{Mumbai, India}

% =======================================================================
\section{Experience}

\entry{Eyecandy Robotics}{Nov 2025 -- Present}
      {Electronics / Hardware Engineer}{}
\begin{points}
  \item Design the PCB stack for a legged companion robot: compute boards (Raspberry Pi Zero 2W, Milk-V SoC), 1S/2S power and battery-charging boards (BQ25887, MP2672), a servo driver board, a mic board and inter-board connectors --- schematic through layout, DFM and fabrication output.
  \item Integrated IMU, buffered servo bus, per-leg force-sensor conditioning, audio amp and mic array on-board, plus test points and jumper options for bring-up. Also wrote the tooling that provisions one runtime image onto a whole fleet, each robot with its own identity.
\end{points}

\entry{NishCorp Technologies}{May -- Nov 2025}
      {Robotics Control Systems Intern}{Hybrid}
\begin{points}
  \item Turned the CAD of a stair-climbing wheelchair into a working simulation --- CAD to URDF, Gazebo/RViz, differential-drive kinematics, IMU plugin, teleop --- so its motion could be studied before it was built in steel. This formed part of the case that helped the startup raise Rs 13 lakh in grants.
  \item Built the half-scale hardware prototype through 3D printing, assembly, motor electronics and field testing; it could sit and stand under its own power.
\end{points}

\entry{Krishna Defence and Allied Industries Ltd.}{May -- Jul 2025}
      {Research Intern, Defence Technology}{at IIT Gandhinagar}
\begin{points}
  \item Solved indoor localization for an autonomous forklift in a naval dockyard, where GPS is unusable and wheel odometry drifts. Chose UWB over Bluetooth, LoRa, RFID, radar and vision, then wrote custom double-sided two-way ranging firmware (Qorvo DWM3001CDK) --- tuning channel, preamble, PAC size and timing delays to take usable range from the stock 45 m to \textbf{315 m}.
  \item Ran SLAM and NAV2 on a Raspberry Pi 4B fusing UWB, IMU and LIDAR; scaled testing from a $2\times2$ m room to a $50\times60$ m yard and validated against a surveyor's total station: \textbf{0.4\% error over 10,000 m\textsuperscript{2}}.
\end{points}

% =======================================================================
\section{Projects}

\pentry{PCB Axial-Flux BLDC Motor}{KiCAD, FEMM, Python, STM32 --- final year project}{2026 -- Present}
\begin{points}
  \item Designing a coreless motor whose stator \emph{is} the PCB: 100 mm 2-layer board, 6 programmatically generated wedge coils over 8 poles, dual 3D-printed rotor, 16 N42 magnets, 0.5 mm air gap.
  \item Built the magnetic model myself: a top-down 2D cut cannot yield torque, so I unrolled it into radial slices to put coil current out of plane, then pulled Kt, Kv and ripple from the air-gap flux. It caught a design error pre-fabrication --- the plastic rotor gives away $3.3\times$ the torque of the steel-backed design it was sized for.
  \item Driver firmware on an STM32G431 + DRV8311H: 20 kHz 3-phase PWM, 40 kHz control ISR, AS5600 encoder over I2C, DMA phase current sensing.
\end{points}

\pentry{Smart Switch Board}{ESP32, ESP-IDF, FreeRTOS, MQTT/TLS, BLE, Flutter, CAD}{2025 -- Present}
\begin{points}
  \item Built an internet-controlled mains wall switch across every layer alone --- CAD enclosure, a relay and optocoupler board isolating 220 V from a 3.3 V MCU, firmware, a TLS MQTT broker, a Flutter app --- under a hard rule that the physical switches keep working with WiFi, cloud and app down. Installed and running at home.
\end{points}

\pentry{e-Yantra Warehouse Drone, IIT Bombay}{ROS2, OpenCV, ArUco, PID}{Oct 2024 -- Jan 2025}
\begin{points}
  \item Tuned PID control for stable hover and precise motion, and wrote the vision that decodes ArUco markers for localization and pick-and-place in a simulated warehouse. \textbf{National semifinalist.}
\end{points}

\pentry{MARIO --- 3-DOF Teaching Arm}{ROS2, Gazebo Fortress, micro-ROS, ESP32, C++}{Jul -- Aug 2024}
\begin{points}
  \item Ported my robotics society's teaching arm off end-of-life Gazebo Classic, rewrote teleop for beginners, and unified the ESP32/micro-ROS pipeline so one command set drives both sim and hardware. Still runs the ROS2 workshops.
\end{points}

\pentry{Also built}{vision/ML, embedded, mechanical}{2023 -- 2025}
\begin{points}
  \item \textbf{Coin sorter} (24 h hackathon --- CNN on a Pi routes coins by servo gate, above 95\% on Rs 5 and Rs 10), \textbf{voice-controlled pick-and-place arm} (speech to detection to inverse kinematics), \textbf{Kurma} quadruped (statically stable turtle gait), \textbf{self-balancing robot}, \textbf{maze solver} (4th), \textbf{gesture-controlled car} (glove IMU over ESP-NOW), a cat feeder and a river-cleaning boat.
\end{points}

% =======================================================================
\section{Leadership \& Teaching}
\begin{points}
  \item \textbf{Society of Robotics and Automation, VJTI} --- Active Member, Aug 2024 -- Mar 2025. Taught workshops on embedded C, ESP32 and sensor fusion and mentored juniors through their first robots; the society builds the kits these run on, and \textbf{over 200 students} have learned on them.
  \item \textbf{Pratibimb, VJTI cultural festival} --- Department Head, Electrical Engineering, Aug 2023 -- Mar 2025. Led the branch to its first overall win in the fest's 26-year history. Also directed a short film, \textit{Plot Se Panga}, which won first prize across all branches.
\end{points}

% =======================================================================
\section{Technical Skills}
\skillrow{Languages}{C, C++, Python, Dart, Lua, TypeScript}
\skillrow{Robotics, control \& embedded}{ROS/ROS2, Gazebo, RViz, URDF, micro-ROS, SLAM, NAV2, PID, sensor fusion, kinematics, FOC; ESP32 (ESP-IDF, FreeRTOS, ESP-NOW), STM32, Raspberry Pi, I2C/SPI/UART, PWM, ADC/DMA, BLE, MQTT}
\skillrow{Electronics, simulation \& tools}{KiCAD schematic \& layout, DFM, buck regulators, battery charging, motor drivers, mains isolation, scope and logic-analyzer bring-up; FEMM magnetostatics, OnShape, 3D printing, OpenCV, ArUco, TensorFlow, Linux, Git, Astro}

\end{document}
